\section{Introduction}\label{sec:intro}

A face recognition system is a distance function wearing a decision rule. An image becomes a vector; two vectors are compared; a threshold turns the comparison into ``same person'' or ``different person''. Everything follows from the map: if one identity's vectors land close together and different identities land far apart, one threshold can work; if not, no threshold tuning will save it.

There are two canonical maps. Eigenfaces applies PCA to centred face images \citep{turk1991}: unsupervised, cheap, interpretable. The modern approach trains a deep network on a large corpus with a metric-learning objective---FaceNet's triplet loss \citep{schroff2015}, VGGFace2's 3.31 million images \citep{cao2018}, ArcFace's angular margin \citep{deng2019}---and uses it as a fixed feature extractor; NIST's FRVT evaluates such systems at scale \citep{grother2019}. Does the learned embedding separate identity better, and can the difference be shown with a number rather than a nice-looking picture? That question is usually answered either by a benchmark table that hides the mechanism or by two scatter plots that are an opinion. We wanted a comparison small enough to inspect completely and controlled enough that the difference cannot be attributed to plotting choices.

\paragraph{Framing.} This is not a novel-method paper; it is a controlled replication and measurement study with a practitioner's case study attached. The comparison is not ``learned versus raw'' in any clean sense: it is a pretrained embedding trained on millions of faces versus full-dimensional raw pixels from 400 Olivetti images. PCA sees the probe images when fitting the \emph{plots}, but the headline scores use unprojected vectors; that transductive visualisation does not make the measured gap conservative. Upsampling $64\times64$ to $160\times160$ invents no detail for Map B, and its training corpus dwarfs the test set. The honest claim is conditional: under these conditions the pretrained embedding separates identity better, and the gap can be measured.

\paragraph{What was done.} Three studies share one codebase and one measurement discipline. (i) \emph{Two maps}: Olivetti's 400 greyscale images of 40 people \citep{samaria1994}, processed as raw pixels (Map A) or frozen embeddings (Map B), with identical PCA settings for visualisation and metrics in full dimensionality. (ii) \emph{One known-population (closed-set) verification protocol for both}: five enrol and five probe per identity, five seeded splits, reporting EER, TAR at fixed FAR, and FAR at matched 1\% FRR. (iii) \emph{A webcam illumination field study}: twelve scans of one person, ten frames each, under six lighting conditions---a case study of whether ``robust to lighting'' holds uniformly, and it does not.

\paragraph{Measurement discipline and findings.} Two-dimensional PCA maps are for human eyes; a projection can look convincing while discarding the geometry a classifier uses, and our raw-pixel map preserves more of its true pairwise structure than the embedding's even though the embedding is better by every metric. We quantify the loss with a rank correlation between 2-D and full-dimensional pairwise distances and treat any statement not backed by a full-dimensional number as a hypothesis. The embedding wins every identity metric, and the margin grows as the task becomes more like a real decision: silhouette 0.106$\to$0.497, 1-NN 94.5\%$\to$100\%, pair ROC-AUC 0.881$\to$0.9995, known-population verification EER 8.7\%$\to$0.22\%. For this one subject, the observed front-lit room-versus-flash contrast falls within repeat-scan variation while flash-direction labels accompany mean distances 3.83 times within-scan jitter; the capture is not randomised, so these are not isolated causal effects. The cross-domain distance to Olivetti is large and domain-confounded; we report it only as a descriptive distance, not as an impostor distribution and not as a bound on stranger distances.

\paragraph{What is not claimed.} No state-of-the-art performance, new embedding, or calibrated security system is claimed: the verification data are small and pre-aligned from a single session; the webcam study has one subject and no stranger; and the verification protocol uses only the 40 known identities, so it does not measure rejection of unknown people.
